%!TEX root = ../thesis.tex
%*******************************************************************************
%*********************************** First Chapter *****************************
%*******************************************************************************

\chapter{INTRODUCTION}  %Title of the First Chapter
\label{chap:intro}

\ifpdf
    \graphicspath{{Chapter1/Figs/Raster/}{Chapter1/Figs/PDF/}{Chapter1/Figs/}}
\else
    \graphicspath{{Chapter1/Figs/Vector/}{Chapter1/Figs/}}
\fi

% ------------------------------ Abbreviations ---------------------------------
% Entries are collected into the Nomenclature list (group Z = acronyms).
\nomenclature[z-VQA]{VQA}{Visual Question Answering}
\nomenclature[z-VLM]{VLM}{Vision-Language Model}
\nomenclature[z-OCR]{OCR}{Optical Character Recognition}
\nomenclature[z-PEFT]{PEFT}{Parameter-Efficient Fine-Tuning}
\nomenclature[z-LoRA]{LoRA}{Low-Rank Adaptation}
\nomenclature[z-QLoRA]{QLoRA}{Quantized Low-Rank Adaptation (LoRA on a 4-bit NF4 base model)}
\nomenclature[z-NF4]{NF4}{4-bit NormalFloat quantization}
\nomenclature[z-EM]{EM}{Exact Match}
\nomenclature[z-VQA-Acc]{VQA-Acc}{VQA Accuracy}
\nomenclature[z-ANLS]{ANLS}{Average Normalized Levenshtein Similarity}
\nomenclature[z-POPE]{POPE}{Polling-based Object Probing Evaluation}
\nomenclature[z-SOTA]{SOTA}{State of the art}

%********************************** Section 1.1 ********************************
\section{Background and Motivation}
\label{sec:intro-background}

\subsection{Visual question answering and vision-language models}

Visual Question Answering (VQA) asks a system to answer a natural-language
question about the content of an image~\cite{antol2015vqa}. Unlike image
classification or object detection, VQA requires joint understanding of the
image, the meaning of the question and, in many cases, the text that appears
in the image~\cite{singh2019textvqa,biten2019stvqa}. Early VQA systems combined
a convolutional image encoder with a recurrent question encoder and selected the
answer from a fixed vocabulary. Recent vision-language models (VLMs) such as
BLIP-2~\cite{li2023blip2}, LLaVA~\cite{liu2023llava},
Florence-2~\cite{xiao2024florence2} and Qwen2.5-VL~\cite{bai2025qwen25vl} take a
different approach: a vision encoder turns the image into a sequence of tokens,
and a large language model reads these tokens together with the question and
generates a free-form answer. These models answer open questions and read text
in images without task-specific training.

\subsection{Vietnamese and English-centric VLMs}

Most open-source VLMs are pretrained mainly on English and Chinese data. When
they are applied to Vietnamese, we observe four recurring problems, labelled
V1--V4 and used throughout this thesis:

\begin{description}
  \item[V1, reading accented text in images.] Vietnamese tone marks are small
  strokes, and misreading one changes the meaning of a word
  (\vi{``BÁNH MÌ''} read as ``banh mi'').
  \item[V2, wrong language or wrong format.] The model answers in English,
  inserts Chinese characters, or returns a full sentence when the reference
  answer is a few words long.
  \item[V3, reading order and word boundaries.] Spaces in Vietnamese separate
  syllables, not words; reading text regions in the wrong order breaks words
  (\vi{``cửa hàng điện thoại''} becomes \vi{``cửa điện hàng thoại''}).
  \item[V4, hallucination.] The answer is fluent but refers to objects or text
  that are not in the image.
\end{description}

The gap between machines and humans on Vietnamese data is still large. On
ViTextVQA~\cite{nguyen2024vitextvqa}, the best published result before this work
is an exact match (EM) of 25.48, while human performance is 88.00.

\subsection{Resource constraints}

Full fine-tuning of a VLM with several billion parameters requires multiple
large GPUs, which is beyond the means of a student team. The university does
not provide a GPU server for this project, so all experiments run on Google
Colab. We therefore make efficiency an explicit design constraint: we use only
parameter-efficient fine-tuning (LoRA~\cite{hu2022lora} and
QLoRA~\cite{dettmers2023qlora}), and every result is reported together with its
training time, inference time and GPU memory.

%********************************** Section 1.2 ********************************
\section{Problem Statement}
\label{sec:intro-problem}

Given an image $I$ and a Vietnamese question $q$, the system generates a
Vietnamese answer $\hat{a} = f(I, q)$. The answer is compared with the reference
answer $a$ after the same text normalization step. The two datasets used in
this thesis pose different requirements:

\begin{itemize}
  \item \textbf{ViVQA}~\cite{tran2021vivqa}: general questions about everyday
  images (colour, counting, objects, position). 99.83\% of the reference answers
  have at most three words, so EM and VQA accuracy are suitable measures.
  \item \textbf{ViTextVQA}~\cite{nguyen2024vitextvqa}: the answer is text read
  from the image, with an average length of 4.18 words. We also report
  ANLS~\cite{biten2019stvqa}, which tolerates small character-level reading
  errors.
\end{itemize}

For example, the ViVQA question \vi{``màu của xe đạp là gì?''} (``what colour
is the bicycle?'') has the answer \vi{``màu đỏ''} (``red''). The ViTextVQA
question \vi{``cửa hàng có biển màu vàng là nơi nào?''} (``what is the shop with
the yellow sign?'') has the answer \vi{``đại lý xổ số miền bắc''} (``northern
lottery agency''), copied directly from the sign.

%********************************** Section 1.3 ********************************
\section{Research Gaps}
\label{sec:intro-gaps}

From the literature review in Chapter~\ref{chap:related}, we identify four gaps:

\begin{description}
  \item[G1. English-centric pretraining.] Zero-shot models answer Vietnamese
  questions partly in English and in the wrong format.
  \item[G2. Weak Vietnamese VQA, especially for scene text.] The state of the art
  on ViTextVQA is 25.48 EM, compared with 88.00 for humans.
  \item[G3. PEFT for Vietnamese VQA with VLMs is unexplored.] The Vietnamese VQA
  studies we reviewed either train task-specific models or use frozen models;
  none applies LoRA or QLoRA to a large VLM.
  \item[G4. No hallucination benchmark for Vietnamese VQA.] There is no
  POPE-style~\cite{li2023pope} benchmark in Vietnamese.
\end{description}

This thesis focuses on G3 and uses G2 and G4 as the settings in which the
proposed method is tested.

%********************************** Section 1.4 ********************************
\section{Objectives}
\label{sec:intro-objectives}

The overall objective is to adapt an open-source VLM to Vietnamese VQA under a
strict parameter budget, and to verify experimentally at least one improvement
over a reproduced baseline. Table~\ref{tab:objectives} breaks this objective
into measurable sub-objectives.

\begin{table}[htbp]
\caption{Objectives and completion criteria}
\label{tab:objectives}
\centering
\small
\begin{tabular}{p{0.30\textwidth}p{0.62\textwidth}}
\toprule
\textbf{Objective} & \textbf{Completion criterion} \\
\midrule
Reproduce exactly one baseline &
Re-run Qwen2.5-VL-3B-Instruct on a benchmark with a published score and obtain a
small difference (TextVQA val: 79.60 versus 79.3 published) \\
Measure the zero-shot baseline on Vietnamese data &
EM, VQA-Acc and ANLS on the full ViVQA and ViTextVQA test sets under one
protocol \\
Adapt with PEFT only &
4-bit QLoRA, fewer than 1\% trainable parameters, one Colab GPU \\
Propose and verify at least one improvement &
Improvement over the baseline under identical conditions, with a paired
statistical test and ablations \\
Analyse errors and hallucination &
Errors broken down by question type and error type; POPE-style hallucination
measurement in Vietnamese \\
Web prototype &
Upload an image, ask in Vietnamese, receive the answer with OCR text and a
confidence score; same answers as the evaluation pipeline \\
Reproducibility &
Every run logs split, seed, hardware, hyperparameters and time; a package with
code, adapters, configurations and an inference script \\
\bottomrule
\end{tabular}
\end{table}

%********************************** Section 1.5 ********************************
\section{Research Questions and Hypotheses}
\label{sec:intro-rq}

\begin{description}
  \item[RQ1.] What does zero-shot inference achieve on ViVQA and ViTextVQA, and
  which Vietnamese-specific errors dominate?
  \item[RQ2.] How much does QLoRA improve over zero-shot inference, without full
  fine-tuning?
  \item[RQ3.] Does accented OCR text in reading order help the model read text in
  images, and does it interact with QLoRA?
  \item[RQ4.] After QLoRA adaptation, does the model hallucinate more or less?
\end{description}

Each hypothesis in Table~\ref{tab:hypotheses} has a quantitative threshold that
was recorded before the corresponding experiment was run.

\begin{table}[htbp]
\caption{Hypotheses and the research questions they answer}
\label{tab:hypotheses}
\centering
\small
\begin{tabular}{p{0.06\textwidth}p{0.66\textwidth}p{0.18\textwidth}}
\toprule
\textbf{ID} & \textbf{Hypothesis} & \textbf{Answers} \\
\midrule
H1 & QLoRA on Qwen2.5-VL-3B increases EM by at least 5 points over the same model
in zero-shot mode. & RQ2; V1--V3 \\
H2 & Accented OCR text ordered from top-left to bottom-right, added to the
prompt, increases ANLS on ViTextVQA by at least 3 points; combined with QLoRA it
gives the best result. & RQ3; V1, V3 \\
H3 & A Vietnamese format prompt (short, accented answer) increases EM by at
least 3 points and reduces wrong-language or unaccented answers. & RQ1; V2 \\
H4 & QLoRA does not increase the hallucination rate compared with zero-shot
inference. & RQ4; V4 \\
\bottomrule
\end{tabular}
\end{table}

%********************************** Section 1.6 ********************************
\section{Contributions}
\label{sec:intro-contrib}

We use Qwen2.5-VL-3B-Instruct~\cite{bai2025qwen25vl} as the baseline, reproduce
it first, and then adapt it to Vietnamese with QLoRA only. The 3-billion-parameter
base model is quantized to 4 bits and kept frozen; we train 7.37 million adapter
parameters, 0.196\% of all parameters, on a single NVIDIA L4 GPU. The
contributions of this thesis are:

\begin{enumerate}
  \item \textbf{OCR-aware QLoRA.} Text read from the image by VietOCR (accents
  preserved, sorted in reading order) is inserted into the prompt during
  training, so that the adapter learns to use it. On the full ViTextVQA test set
  (10,028 questions), this reaches 53.35 EM versus 50.26 for plain QLoRA on the
  same data ($+3.09$ EM, McNemar test $p = 6.7\times10^{-21}$), above the best
  published result on the same test set (ViTextBLIP-2, 25.48
  EM~\cite{nguyen2024vitextvqa}).
  \item \textbf{A separated mechanism.} A $2\times2$ experiment shows that OCR
  helps only when it is used both during training and at inference: adding OCR
  to the prompt of a model that never saw OCR does not improve results.
  \item \textbf{A strong Vietnamese general-VQA result.} On ViVQA, QLoRA raises
  EM from 22.09 (zero-shot) to 76.64.
  \item \textbf{A Vietnamese hallucination benchmark.} POPE-vi, a POPE-style
  benchmark with 4,985 yes/no questions on 500 images not seen during training.
  \item \textbf{Analysis and a verified prototype.} Ablations for each component,
  automatic and manual error analysis, and a web prototype verified to produce
  the same answers as the official evaluation pipeline.
\end{enumerate}

All conclusions rest on decision rules recorded before the experiments, paired
tests on the same test set, and complete logs (split, seed, GPU, time), so every
number can be re-run from the public repository.

%********************************** Section 1.7 ********************************
\section{Scope and Limitations}
\label{sec:intro-scope}

\begin{table}[htbp]
\caption{In scope and out of scope}
\label{tab:scope}
\centering
\small
\begin{tabular}{p{0.45\textwidth}p{0.45\textwidth}}
\toprule
\textbf{In scope} & \textbf{Out of scope} \\
\midrule
ViVQA and ViTextVQA with fixed splits &
Building new training datasets or data augmentation \\
One baseline, Qwen2.5-VL-3B-Instruct; other VLMs for reference only &
Comparing several backbones as main baselines \\
Adaptation with PEFT only (4-bit QLoRA); OCR text as an auxiliary input &
Full fine-tuning; training a VLM from scratch \\
EM, VQA-Acc, ANLS (token F1 as a secondary metric); ablations, error and
hallucination analysis &
Retrieval-augmented VQA with external knowledge \\
Web prototype running on Colab with one GPU &
Large-scale deployment or mobile applications \\
\bottomrule
\end{tabular}
\end{table}

\paragraph{Assumptions.} We assume that the official splits of both datasets
represent the question distribution of each dataset, that the reference answers
are correct for the large majority of questions (incorrect references are
measured and reported separately in Chapter~\ref{chap:results}), and that an L4
GPU on Colab gives stable results across sessions when the seed and
configuration are fixed.

\paragraph{Known limitations.} Results are bounded by a 3-billion-parameter
base model at 4 bits and by OCR quality. Main experiments use one seed; the OCR
comparison is additionally checked on three seeds. EM is strict towards answers
that are correct in meaning but written differently, so we add ANLS and manual
review. These limitations are revisited in Chapters~\ref{chap:discussion}
and~\ref{chap:conclusion}.

%********************************** Section 1.8 ********************************
\section{Thesis Outline}
\label{sec:intro-outline}

Chapter~\ref{chap:related} reviews VQA, vision-language models,
parameter-efficient fine-tuning and prior work on Vietnamese VQA.
Chapter~\ref{chap:methods} describes the datasets, the baseline, the proposed
OCR-aware QLoRA method and the evaluation protocol. Chapter~\ref{chap:results}
reports the main results and ablations. Chapter~\ref{chap:discussion} discusses
error analysis, hallucination and limitations. Chapter~\ref{chap:conclusion}
concludes the thesis. Appendix~\ref{appendix:pm} presents the project management
plan.
